\lyseis{7.5}{Πολλαπλασιασμός ρητών}
\sol{1} $-12,\ -8,\ -4,\ 0,\ 4,\ 8,\ 12$.
\sol{2} α) $-24$ \quad β) $-21$ \quad γ) $45$ \quad δ) $0$ \quad ε) $40$ \quad στ) $1$.
\sol{3} α) $-6$ \quad β) $3{,}6$ \quad γ) $-1$ \quad δ) $\frac{2}{3}$ \quad ε) $-6$ \quad στ) $\frac{1}{3}\cdot(-3)=-1$.
\sol{4} α) Αρνητικό ($3$ αρνητικοί). β) Θετικό ($4$). γ) Αρνητικό ($3$). δ) Είναι $0$.
\sol{5} α) $-8$ \quad β) $8$ \quad γ) $0$ \quad δ) $7$.
\sol{6} α) $-40$ \quad β) $25$ \quad γ) $-10+18=8$ \quad δ) $-594$ \quad ε) $-4-3=-7$.
\sol{7} α) $-12+10+6=4$ \quad β) $(-2)(-3)+3(-3)=6-9=-3$ \quad γ) $10-30=-20$.
\sol{8} α) Λάθος: $+12$. β) Σωστό. γ) Σωστό. δ) Σωστό. ε) Λάθος: αρνητικό. στ) Σωστό.
\sol{9} Σε $3$ ώρες $3\cdot(-2)=-6$ βαθμούς. Πριν $3$ ώρες $(-2)\cdot(-3)=+6$, δηλαδή τότε είχε $6$ βαθμούς περισσότερο.
\sol{10} $5\cdot(-4)=-20$ €.
\sol{11} Σε $8$ λεπτά: $8\cdot(-3)=-24$ μ. Πριν $5$ λεπτά: $(-3)\cdot(-5)=+15$, δηλαδή $15$ μ. πιο ψηλά.
\sol{12} $4\cdot(-2)=-8$ όροφοι (κατεβαίνει $8$).
\sol{13} $9\cdot(-4)=-36$ €.
\sol{14} $6\cdot(-0{,}25)=-1{,}5$ €.
\sol{15} Αν η θερμοκρασία πέφτει $2$ βαθμούς την ώρα, πηγαίνοντας $3$ ώρες πίσω έχουμε $6$ περισσότερους. Στο μοτίβο: $0\cdot(-3)=0$, $(-1)\cdot(-3)=3$, $(-2)\cdot(-3)=6$.
\sol{16} Γιατί το $\alpha\cdot\alpha$ έχει ομόσημους παράγοντες (ίδιο πρόσημο), άρα θετικό (ή $0$).
\sol{17} Άρτιο πλήθος αρνητικών, π.χ. $2$, $4$, $6$, $8$, $10$ αρνητικοί.
\sol{18} Ένας τουλάχιστον από τους δύο είναι $0$.
\sol{19} Με $100$ παράγοντες $+1$ και με $101$ παράγοντες $-1$.
\sol{20} Ελεύθερη απάντηση: $3\cdot(-5)=-15$, $2\cdot(-5)=-10$, \dots, $(-3)\cdot(-5)=15$.
